\section{行為準則}

\subsection{事業主等の行為準則}

\begin{InyouJobunbox}
  \begin{itemize}
    \item 法第69条〜第72条
    \item 則第85条の2〜第86条
  \end{itemize}
\end{InyouJobunbox}

\begin{fillblank}[自作 2026/11/13]
  \noindent \textbf{◯DB則}
  \begin{lawstructure}
    \begin{lawarticle}{第八十五条の三}{（　　A　　）に係る事業主への情報提供}
      第八条第二項の代表は、（　　A　　）をしようとするときは、当該変更に係る実施事業所の事業主（当該代表を除く。）に対し、遅滞なく、（　　B　　）に関する情報の提供を行わなければならない。

      \lawparagraph{2}
      基金は、（　　A　　）をしようとするときは、当該変更に係る実施事業所の事業主に対し、遅滞なく、（　　B　　）に関する情報の提供を行わなければならない。

    \end{lawarticle}
  \end{lawstructure}
\end{fillblank}

\begin{answerbox}
  \begin{enumerate}[A: ]
    \item 規約の変更
    \item 変更の内容及び規約変更日
  \end{enumerate}
\end{answerbox}

\newpage

\subsection{業務概況の周知}

\begin{InyouJobunbox}
  \begin{itemize}
    \item 法第73条
    \item 則第87条
    \item 解釈第8-3〜8-4
  \end{itemize}
\end{InyouJobunbox}

\begin{fillblank}[年金1 2024 問題1(1)]
  \noindent \textbf{◯DB則}
  \begin{lawstructure}
    \begin{lawarticle}{第八十七条}{業務概況の周知}
      事業主等（第七号に掲げる事項については第八十二条の要件に該当する規約型企業年金の事業主を除き、第八号に掲げる事項についてはリスク分担型企業年金を実施する事業主等に限る。）が法第七十三条第一項の規定に基づき、その確定給付企業年金に係る業務の概況について加入者に周知させる場合においては、毎事業年度一回以上、当該時点における次に掲げる事項（第二号から第六号までに掲げる事項にあっては、当該時点における直近の概況。以下この条において「周知事項」という。）を加入者に周知させるものとする。

      \lawitem{一}
      給付の種類ごとの標準的な給付の額及び給付の設計

      \lawitem{二}
      加入者の数及び給付の種類ごとの受給権者の数

      \lawitem{三}
      給付の種類ごとの給付の支給額その他給付の支給の概況

      \lawitem{四}
      事業主が資産管理運用機関等に納付した掛金の額、納付時期その他掛金の納付の概況

      \lawitem{五}
      積立金の額と（　　A　　）との比較その他積立金の積立ての概況

      \lawitem{六}
      積立金の運用収益又は運用損失及び資産の構成割合その他積立金の運用の概況

      \lawitem{七}
      （　　B　　）の概要

      \lawitem{八}
      調整率の推移その他調整率に関する事項

      \lawitem{九}
      その他確定給付企業年金の事業に係る重要事項

      \lawparagraph{2}
      （略）

      \lawparagraph{3}
      事業主等が加入者に周知事項を周知させる場合であって、前項各号のいずれかの方法を選択するときは、加入者以外の者であって事業主等が（　　C　　）を負っているものにも周知が行われる方法を選択するよう努めなければならない。

      \lawparagraph{4}
      （　　D　　）を実施する事業主等は、前項の規定にかかわらず、毎事業年度一回以上、周知事項を加入者以外の者であって事業主等が（　　C　　）を負っているものに周知させるものとする。
    \end{lawarticle}
  \end{lawstructure}
\end{fillblank}

\begin{answerbox}
  \begin{enumerate}[A: ]
    \item 責任準備金の額及び最低積立基準額
    \item 基本方針
    \item 給付の支給に関する義務
    \item リスク分担型企業年金
  \end{enumerate}
\end{answerbox}

\begin{truefalse}[年金1 2019 問題1(2)\textcircled{2}]
  下線部の正誤を判定し、誤っている場合は正しい内容に改めなさい。
\\

  リスク分担型企業年金を実施する事業主等が毎事業年度一回以上、加入者に周知させる業
務概況のうち「調整率の推移その他調整率に関する事項」は、 「確定給付企業年金法並びにこ
れに基づく政令及び省令について（法令解釈） 」によると、例えば、次に掲げる事項を定める
ことである。
\begin{enumerate}[A. ]
  \item 年金額の改定ルール
  \item \underline{過去3年程度の調整率の推移および今後の見通し}
  \item 「B」の調整率の算出根拠となったデータ
  \item その他、調整率に重要な影響を与えると認められる事項
\end{enumerate}
\end{truefalse}

\begin{answerbox}
  \textbf{誤}

  過去5年程度の調整率の推移
\end{answerbox}

\newpage